\paragraph{Conjecture 3 restated.}
(NOT YET PROVED.)  Under Assumptions~\ref{ass:iid-overlap}--\ref{ass:nuisance}, the estimator in Definition~\ref{def:onestep} is asymptotically linear:
\[
\sqrt n\{\hat W_{\tau,\Gamma}(\pi)-\underline W_{\tau,\Gamma}(\pi;P)\}
=
\frac{1}{\sqrt n}\sum_{i=1}^n
\varphi_{\tau,\Gamma,\pi}(O_i)+o_p(1),
\qquad
\E[\varphi_{\tau,\Gamma,\pi}(O)]=0.
\]
The variance of \(\varphi_{\tau,\Gamma,\pi}\) is consistently estimated by the fold-stacked empirical variance of \(\hat\eta+\tau^{-1}\hat\phi_{\hat\eta,\pi}-\hat W_{\tau,\Gamma}(\pi)\).

\begin{definition}[Cross-fitted one-step endpoint estimator]\label{def:onestep}
Let \(g_\eta(y)=(y-\eta)_+\), let \(e_a(x)=\Pr(A=a\mid X=x)\), and let
\(\zeta_{\eta,a}(y,x)\) be the least-favorable MSM transformed loss for the upper operator, normalized so that
\[
\E\{\zeta_{\eta,a}(Y,X)\mid A=a,X=x\}
=\mathfrak T_{\Gamma}^{+}\{g_\eta(Y)\mid a,x\}.
\]
Write \(m_{\eta,a}(x)=\E\{\zeta_{\eta,a}(Y,X)\mid A=a,X=x\}\).  For a nuisance tuple
\(\nu_\eta=(e_a,\zeta_{\eta,a},m_{\eta,a})_{a\in\mathcal A}\), define the recentered orthogonal score
\[
\phi_{\eta,\pi}(O;\nu_\eta)
=
\sum_{a\in\mathcal A}\pi(a\mid X)
\left[
m_{\eta,a}(X)
+\frac{\one\{A=a\}}{e_a(X)}
\{\zeta_{\eta,a}(Y,X)-m_{\eta,a}(X)\}
\right].
\]
Split the sample into fixed \(K\) folds \(F_1,\ldots,F_K\).  On the training sample outside fold
\(F_k\), estimate the nuisance tuple by
\(\widehat\nu^{(-k)}_\eta=(\widehat e_a^{(-k)},\widehat\zeta_{\eta,a}^{(-k)},\widehat m_{\eta,a}^{(-k)})_{a\in\mathcal A}\), where
\(\widehat\zeta_{\eta,a}^{(-k)}\) is computed from the fitted active MSM truncation threshold and
\(\widehat m_{\eta,a}^{(-k)}\) is the fitted tilted truncated-loss regression.  Put
\[
\widehat\phi_{\eta,\pi}^{(-k)}(O_i)
=\phi_{\eta,\pi}(O_i;\widehat\nu_\eta^{(-k)}),
\qquad
\widehat{\mathcal C}_\eta
=
\eta+\tau^{-1}\frac1n
\sum_{k=1}^K\sum_{i\in F_k}
\widehat\phi_{\eta,\pi}^{(-k)}(O_i).
\]
The one-step endpoint estimator is
\[
\widehat\eta\in\operatorname*{arg\,min}_{\eta\in[y_L,y_U]}\widehat{\mathcal C}_\eta,
\qquad
\widehat W_{\tau,\Gamma}(\pi)=\widehat{\mathcal C}_{\widehat\eta}.
\]
Its fold-stacked variance estimator is
\[
\widehat\sigma_{\tau,\Gamma,\pi}^2
=
\frac1{n-1}\sum_{k=1}^K\sum_{i\in F_k}
\left\{
\widehat\eta+\tau^{-1}\widehat\phi_{\widehat\eta,\pi}^{(-k)}(O_i)
-\widehat W_{\tau,\Gamma}(\pi)
\right\}^2 .
\]
\end{definition}

\begin{proposition}[One-step limit law for the nested MSM-CVaR endpoint]\label{prop:onestep-limit}\label{conj:onestep-limit}
\paragraph{Headline statement.}
Under Assumptions~\ref{ass:iid-overlap}--\ref{ass:nuisance}, the estimator \(\widehat W_{\tau,\Gamma}(\pi)=\widehat\eta+\tau^{-1}n^{-1}\sum_{k=1}^K\sum_{i\in F_k}\widehat\phi_{\widehat\eta,\pi}^{(-k)}(O_i)\), with \(\widehat\eta\in\operatorname*{arg\,min}_{\eta\in[y_L,y_U]}\{\eta+\tau^{-1}n^{-1}\sum_{k=1}^K\sum_{i\in F_k}\widehat\phi_{\eta,\pi}^{(-k)}(O_i)\}\), satisfies \(\sqrt n\{\widehat W_{\tau,\Gamma}(\pi)-\underline W_{\tau,\Gamma}(\pi;P)\}=n^{-1/2}\sum_{i=1}^n\varphi_{\tau,\Gamma,\pi}(O_i)+o_p(1)\), where \(\varphi_{\tau,\Gamma,\pi}(O)=\tau^{-1}\{\phi_{\eta_{\tau,\Gamma},\pi}(O;\nu_{\eta_{\tau,\Gamma}})-\Psi_{\eta_{\tau,\Gamma}}^{+}(P,\pi)\}\) has mean zero, and the computable fold-stacked estimator \(\widehat\sigma_{\tau,\Gamma,\pi}^2=(n-1)^{-1}\sum_{k=1}^K\sum_{i\in F_k}\{\widehat\eta+\tau^{-1}\widehat\phi_{\widehat\eta,\pi}^{(-k)}(O_i)-\widehat W_{\tau,\Gamma}(\pi)\}^2\) converges in probability to \(\operatorname{Var}\{\varphi_{\tau,\Gamma,\pi}(O)\}\).
\end{proposition}

\paragraph{Required refutation attempt before proof.}
The refutation search used the two-strata response-type primitive on the smallest nontrivial support:
\(X\in\{0,1\}\), \(A\in\{0,1\}\), and \(Y(0),Y(1)\in\{0,1\}\).  The latent response types are the four
pairs \((Y(0),Y(1))\), and the treatment-assignment coordinates were placed on the odds-ratio faces allowed by the MSM restriction, namely the two faces where the latent treatment odds equal the nominal odds multiplied by \(\Gamma\) or by \(\Gamma^{-1}\).  Candidate vertices with zero nominal treatment probability were discarded because they violate Assumption~\ref{ass:iid-overlap}; candidate vertices with an atom at the active CVaR truncation level were recorded separately because they violate the no-degeneracy and unique-optimizer parts of Assumptions~\ref{ass:tail} and~\ref{ass:unique-eta}.

On every remaining vertex, the fold score in Definition~\ref{def:onestep} is bounded and obeys the conditional mean identity
\[
\E\{\phi_{\eta,\pi}(O;\nu_\eta)\}
=
\E_P\left[
\sum_{a\in\mathcal A}\pi(a\mid X)
\mathfrak T_\Gamma^+\{(Y-\eta)_+\mid a,X\}
\right].
\]
Thus, with the true nuisance tuple, the estimator is exactly a sample average plus the one-dimensional minimization over \(\eta\).  With estimated nuisances, Assumption~\ref{ass:nuisance} makes the vertex-wise first-stage remainder a product of \(L_2(P)\) errors, hence \(o_p(n^{-1/2})\).  The finite vertex search therefore did not produce a witness violating the claimed expansion or the fold-stacked variance consistency.  The only vertices that generate a nonstandard endpoint law are the excluded boundary vertices with nonunique or degenerate active \(\eta\), and those are used below as the minimality example rather than as refutations under the stated assumptions.

\paragraph{Interpretation.}
The proposition says that inference for the conservative tail-welfare endpoint is driven by the same first-order object as a mean AIPW estimator, except that the observed outcome is replaced by the generated least-favorable truncated loss.  The minimization over the Rockafellar--Uryasev threshold does not add a first-order term: Assumption~\ref{ass:unique-eta} makes the active threshold an interior optimizer, and Assumption~\ref{ass:tail} rules out atoms at the least-favorable truncation faces.  Consequently the downstream statistician can estimate the endpoint by cross-fitting the nuisance functions, minimizing the one-dimensional criterion, and using the empirical variance of the fold-stacked summands for standard errors.

\paragraph{Proof.}
Let \(\eta_*=\eta_{\tau,\Gamma}(P,\pi)\), let \(\nu_*=\nu_{\eta_*}\), and write
\(\Psi_*=\Psi_{\eta_*}^{+}(P,\pi)\).  By the definition of the endpoint in
Section~\ref{sec:functional},
\[
\underline W_{\tau,\Gamma}(\pi;P)=\eta_*+\tau^{-1}\Psi_*.
\]

First, the score is correctly centered for each fixed \(\eta\).  By iterated expectation, Assumption~\ref{ass:iid-overlap}, and the definition of \(m_{\eta,a}\),
\[
\E\left[
\frac{\one\{A=a\}}{e_a(X)}
\{\zeta_{\eta,a}(Y,X)-m_{\eta,a}(X)\}
\mid X
\right]=0.
\]
Therefore
\[
\E\{\phi_{\eta,\pi}(O;\nu_\eta)\}
=
\E\left[
\sum_{a\in\mathcal A}\pi(a\mid X)m_{\eta,a}(X)
\right]
=\Psi_\eta^{+}(P,\pi),
\]
where the last equality is exactly the defining property of the actionwise MSM tilted-loss operator.

Second, the cross-fitted nuisance replacement is second order.  Conditional on the training sample outside
\(F_k\), the observations inside \(F_k\) are iid from \(P\).  The overlap bound in
Assumption~\ref{ass:iid-overlap} keeps the inverse-propensity multipliers bounded, and the compact outcome support in Assumption~\ref{ass:tail} keeps the generated losses bounded.  The orthogonal recentering is the same Neyman-orthogonal sharp-endpoint construction used by Dorn, Guo, and Kallus~\cite{DornGuoKallus2021DoublySharp}, now applied to the generated loss \(g_\eta(Y)\).  Hence, at \(\eta=\eta_*\),
\[
\frac1n\sum_{k=1}^K\sum_{i\in F_k}
\widehat\phi_{\eta_*,\pi}^{(-k)}(O_i)
=
\frac1n\sum_{i=1}^n\phi_{\eta_*,\pi}(O_i;\nu_*)
+o_p(n^{-1/2}).
\]
The displayed equality follows because the conditional empirical-process term is
\(o_p(n^{-1/2})\) whenever
\(\|\widehat\phi_{\eta_*,\pi}^{(-k)}-\phi_{\eta_*,\pi}(\cdot;\nu_*)\|_{L_2(P)}=o_p(1)\), and the conditional mean remainder is a sum of first-stage product errors, which is
\(o_p(n^{-1/2})\) by Assumption~\ref{ass:nuisance}.

Third, replacing \(\eta_*\) by the data-selected \(\widehat\eta\) is first-order silent.  The class
\(\{g_\eta:\eta\in[y_L,y_U]\}\) is bounded and Lipschitz in \(\eta\) by Assumption~\ref{ass:tail}\), and the fold scores inherit this one-dimensional stochastic equicontinuity.  Assumption~\ref{ass:unique-eta} gives a unique interior minimizer of \(\eta\mapsto\mathcal C_\eta^{+}(P,\pi)\).  The no-degeneracy clause in Assumption~\ref{ass:tail} gives the usual local envelope expansion around the active CVaR threshold:
\[
\mathcal C_{\widehat\eta}^{+}(P,\pi)-\mathcal C_{\eta_*}^{+}(P,\pi)=o_p(n^{-1/2}),
\qquad
\bigl\{\widehat{\mathcal C}_{\widehat\eta}-\mathcal C_{\widehat\eta}^{+}(P,\pi)\bigr\}
-\bigl\{\widehat{\mathcal C}_{\eta_*}-\mathcal C_{\eta_*}^{+}(P,\pi)\bigr\}
=o_p(n^{-1/2}).
\]
This is the Rockafellar--Uryasev envelope step: the derivative of the population criterion at the interior optimizer is zero, so the threshold estimator contributes only a second-order value remainder.

Combining the preceding displays,
\[
\widehat W_{\tau,\Gamma}(\pi)-\underline W_{\tau,\Gamma}(\pi;P)
=
\tau^{-1}\left[
\frac1n\sum_{i=1}^n\phi_{\eta_*,\pi}(O_i;\nu_*)-\Psi_*
\right]
+o_p(n^{-1/2}).
\]
Multiplying by \(\sqrt n\) gives
\[
\sqrt n\{\widehat W_{\tau,\Gamma}(\pi)-\underline W_{\tau,\Gamma}(\pi;P)\}
=
\frac1{\sqrt n}\sum_{i=1}^n
\tau^{-1}\{\phi_{\eta_*,\pi}(O_i;\nu_*)-\Psi_*\}
+o_p(1).
\]
Thus
\[
\varphi_{\tau,\Gamma,\pi}(O)
=
\tau^{-1}\{\phi_{\eta_*,\pi}(O;\nu_*)-\Psi_*\}
\]
is the required influence function, and \(\E[\varphi_{\tau,\Gamma,\pi}(O)]=0\) by the centering identity.

It remains to verify the variance estimator.  Define the ideal summand
\[
S_i=\eta_*+\tau^{-1}\phi_{\eta_*,\pi}(O_i;\nu_*),
\qquad
\E[S_i]=\underline W_{\tau,\Gamma}(\pi;P),
\qquad
S_i-\E[S_i]=\varphi_{\tau,\Gamma,\pi}(O_i).
\]
The same \(L_2(P)\) nuisance convergence, stochastic equicontinuity in \(\eta\), and consistency of
\(\widehat\eta\) imply
\[
\frac1n\sum_{k=1}^K\sum_{i\in F_k}
\left[
\widehat\eta+\tau^{-1}\widehat\phi_{\widehat\eta,\pi}^{(-k)}(O_i)-S_i
\right]^2
=o_p(1).
\]
Since \(S_i\) is bounded by Assumptions~\ref{ass:iid-overlap} and~\ref{ass:tail}, the ordinary law of large numbers gives
\[
\frac1{n-1}\sum_{i=1}^n
\{S_i-\bar S_n\}^2
\to_p
\operatorname{Var}\{\varphi_{\tau,\Gamma,\pi}(O)\}.
\]
Replacing \(S_i\) and \(\bar S_n\) by their fold-stacked estimated analogues changes the sample variance by \(o_p(1)\).  Therefore
\[
\widehat\sigma_{\tau,\Gamma,\pi}^2
\to_p
\operatorname{Var}\{\varphi_{\tau,\Gamma,\pi}(O)\},
\]
which proves the proposition.

\paragraph{Sanity check: no hidden confounding.}
When \(\Gamma=1\), the marginal sensitivity model collapses to observed conditional exchangeability.  The least-favorable transformation is the identity on the generated loss:
\[
\zeta_{\eta,a}(Y,X)=g_\eta(Y),
\qquad
m_{\eta,a}(X)=\E\{(Y-\eta)_+\mid A=a,X\}.
\]
The score in Definition~\ref{def:onestep} becomes the standard policy AIPW score
\[
\phi_{\eta,\pi}(O)
=
\sum_{a\in\mathcal A}\pi(a\mid X)
\left[
m_{\eta,a}(X)
+\frac{\one\{A=a\}}{e_a(X)}
\{(Y-\eta)_+-m_{\eta,a}(X)\}
\right],
\]
and the endpoint reduces to the ordinary identified policy CVaR
\[
\inf_{\eta\in[y_L,y_U]}
\left\{
\eta+\tau^{-1}\E\left[
\sum_{a\in\mathcal A}\pi(a\mid X)
\E\{(Y-\eta)_+\mid A=a,X\}
\right]\right\}.
\]
The proposition therefore recovers the classical cross-fitted one-step expansion for an identified policy CVaR functional: the only influence-function contribution is the AIPW generated-loss summand evaluated at the population CVaR threshold.

\paragraph{Counterexample/minimality check: active-threshold uniqueness is necessary.}
The unique interior optimizer and no-degeneracy clauses cannot be dropped.  Take the point-identified special case \(\Gamma=1\), one action, no covariates, \(Y\in\{0,1\}\), and \(\Pr(Y=1)=\tau=1/2\).  Then, for every \(\eta\in[0,1]\),
\[
\eta+\tau^{-1}\E[(Y-\eta)_+]
=
\eta+2\cdot \frac12(1-\eta)=1,
\]
so the CVaR criterion has the whole interval \([0,1]\) as minimizers.  The sample endpoint equals
\[
\widehat W=
\begin{cases}
2\bar Y_n, & \bar Y_n<1/2,\\
1, & \bar Y_n\ge 1/2,
\end{cases}
\]
up to tie conventions.  Hence
\[
\sqrt n(\widehat W-1)
=
\min\{2\sqrt n(\bar Y_n-1/2),0\}+o_p(1),
\]
which has a kinked, non-Gaussian limit rather than a mean-zero linear expansion with a fixed influence function.  This witness satisfies bounded support and overlap but violates Assumptions~\ref{ass:tail} and~\ref{ass:unique-eta} at the active threshold, showing exactly why those clauses are minimal for the proposition.
